%% ============================================================
%%  chapitre9.tex — الفصل التاسع: الاجتهاد القضائي والتجارب المقارنة
%% ============================================================

\chapter{الاجتهاد القضائي والتجارب المقارنة}

%% ─── مقدمة الفصل ──────────────────────────────────────────
\section*{مقدمة الفصل}
\addcontentsline{toc}{section}{مقدمة الفصل}

يحتل الاجتهاد القضائي مكانة محورية في المنظومة القانونية المعاصرة، لا سيما في مادة قانون الشغل التي تتسم بالديناميكية وسرعة التأثر بالتحولات الاقتصادية والاجتماعية. فإذا كانت النصوص التشريعية ترسم الأطر العامة وتحدد القواعد المجردة، فإن القضاء هو الذي يبث فيها الروح ويكتشف إرادة المشرع الحقيقية عند تطبيقها على الوقائع والمنازعات المعروضة أمامه.

وتكتسي دراسة الاجتهاد القضائي في مجالي التشغيل المؤقت والمناولة أهمية مضاعفة، بالنظر إلى ما يشهده الواقع العملي من محاولات متكررة للالتفاف على مقتضيات القانون والتحايل على مبدأ استقرار العمل عبر توظيف هذين النمطين في غير مقاصدهما المشروعة. من هذا المنطلق، نهض القضاء الاجتماعي المغربي—وعلى رأسه محكمة النقض (المجلس الأعلى سابقاً)—بدور حمائي بالغ الأهمية لإرساء التوازن العقدي وحماية الطرف الضعيف في العلاقة الشغلية.

بالموازاة مع ذلك، يقتضي التحليل القانوني الرصين الانفتاح على التجارب المقارنة والمعايير الدولية، للاطلاع على كيفية تعاطي الأنظمة القانونية الرائدة مع إشكالية التوفيق بين متطلبات المرونة الاقتصادية للمقاولات وضمانات الأمان الوظيفي والاجتماعي للأجراء.

يتناول هذا الفصل بالدراسة والتحليل المحاور الأربعة المحددة في خطة الدراسة:
\begin{itemize}[leftmargin=*, itemsep=4pt]
    \item أولاً: الاجتهاد القضائي المغربي (المبادئ والقرارات المرجعية).
    \item ثانياً: موقف القضاء من المناولة المقنعة وتوريد اليد العاملة غير المشروع.
    \item ثالثاً: التجربة المقارنة (الفرنسية، والإسبانية، والألمانية، والتجارب العربية).
    \item رابعاً: المعايير الدولية للشغل ومبادئ منظمة العمل الدولية.
\end{itemize}

%% ─────────────────────────────────────────────────────────
%%  أولاً: الاجتهاد القضائي المغربي
%% ─────────────────────────────────────────────────────────
\section{الاجتهاد القضائي المغربي}

استقر قضاء محكمة النقض (المجلس الأعلى سابقاً) ومحاكم الموضوع بالمغرب على جملة من المبادئ الحمائية الرامية إلى صيانة استقرار الشغل ومنع جعل المرونة التعاقدية مدخلاً لإهدار حقوق الطبقة العاملة.

\subsection{مبدأ سمو الواقع التنفيذي على ظاهر العقد (تكييف العقود)}
يُعد مبدأ «الواقع يسمو على الظاهر» من الركائز الأساسية التي يعتمدها القاضي الاجتماعي؛ إذ تملك المحكمة كامل السلطة في إضفاء التكييف القانوني الصحيح على العلاقة الرابطة بين الطرفين وفقاً لحقيقة تنفيذها الميداني، دون التقيد بالتسمية أو الوصف الشكلي الذي يمنحه المشغل للعقد (تطبيقاً للفصل \textenglish{461} من قانون الالتزامات والعقود والمبادئ العامة لمدونة الشغل).

فإذا أطلقت المقاولة على العقد وصف «عقد مناولة» أو «مهمة تشغيل مؤقت»، في حين كشفت الوقائع عن إدماج الأجير في الهيكل التنظيمي الدائم للمقاولة المستعملة وممارسته لمهام قارة ومستمرة، فإن القضاء يجرد العقد من وصفه الظاهري ويعيد تكييفه باعتباره عقد شغل أصلياً وغير محدد المدة (\textenglish{CDI}) يربط الأجير بالمقاولة المستعملة مباشرة.

\subsection{معيار التبعية القانونية وتحديد المشغل الحقيقي}
تعتبر المادة \textenglish{6} من مدونة الشغل أن عنصر التبعية هو الفيصل الجوهري في قيام عقد الشغل. وقد كرس الاجتهاد القضائي معيار التبعية القانونية—المتمثل في خضوع الأجير لإشراف وتوجيه ورقابة المشغل وسلطته التأديبية—لتحديد المسؤول الحقيقي عن حقوق الأجير:
\begin{itemize}[leftmargin=*, itemsep=6pt]
    \item \textbf{في التشغيل المؤقت القانوني:} تبقى التبعية القانونية والالتزامات التعاقدية من أداء للأجور والتصريح لدى الصندوق الوطني للضمان الاجتماعي ملقاة على عاتق مقاولة التشغيل المؤقت، في حين تقتصر المقاولة المستعملة على ممارسة سلطة توجيهية وإشرافية ظرفية طيلة مدة المهمة المحددة قانوناً.
    \item \textbf{في حالة الانحراف والتحايل:} إذا تجاوزت المقاولة المستعملة حدود الإشراف المؤقت وأخضعت الأجير لسلطتها الرئاسية الدائمة والتأديبية واستخدمته في صلب نشاطها الأساسي، فإن القضاء يعتبر المقاولة المستعملة هي المشغل الفعلي والقانوني، وتتحمل بالتالي التبعات المترتبة عن إنهاء الخدمة.
\end{itemize}

\subsection{مبدأ المساواة في الأجر بين الأجير المؤقت والأجير الدائم}
يُعد مبدأ المساواة في الأجر والامتيازات بين الأجير المؤقت والأجير الدائم من المبادئ الجوهرية التي كرسها المشرع المغربي بمقتضى المادة \textenglish{504} من مدونة الشغل، التي تنص صراحة على أنه \textit{«لا يجوز أن يقل الأجر الذي يتقاضاه الأجير المؤقت عن الأجر المؤدى للأجير الذي يشغل نفس المنصب في المقاولة المستعملة ويتمتع بنفس المؤهلات المهنية.»}

وقد ذهب القضاء المغربي في تأويله لهذا النص إلى إضفاء طابع آمر وغير قابل للتنازل على هذا المبدأ، بحيث لا يمكن لاتفاق الطرفين أو لعقد الوضع رهن الإشارة أن يتضمن شروطاً تُنقص من حق الأجير المؤقت في المساواة. ويشمل مفهوم المساواة في هذا السياق ليس فقط الأجر الأساسي، بل كذلك التعويضات والمنح والامتيازات العينية المرتبطة بالمنصب، كالتعويض عن التنقل، وعلاوة الأقدمية المحتسبة على أساس مدة المهمة، والمنح المتعلقة بظروف العمل الخاصة كالعمل الليلي أو في بيئات خطرة.

وتجد هذه القاعدة أساسها الفلسفي في مبدأ «الأجر المتساوي للعمل المتساوي القيمة» المكرس في الاتفاقية الدولية رقم \textenglish{100} لمنظمة العمل الدولية، والذي يرمي إلى القضاء على كل أشكال التمييز غير المبرر في الاستحقاقات المالية بين العمال الذين يؤدون نفس المهام في نفس ظروف العمل، بصرف النظر عن طبيعة رابطتهم التعاقدية.

\subsection{مبدأ حظر التعسف في إنهاء مهمة الأجير المؤقت}
أرسى الاجتهاد القضائي المغربي مبدأ مهماً مفاده أن إنهاء مهمة الأجير المؤقت قبل الأجل المتفق عليه في عقد الوضع رهن الإشارة، دون وجود سبب حقيقي وجدي أو خطأ جسيم يبرره، يُشكل إنهاءً تعسفياً يُرتب حق الأجير في التعويض. فالمقاولة المستعملة التي تُنهي مهمة الأجير قبل أوانها بقرار أحادي الجانب دون مبرر مشروع، تتحمل المسؤولية التضامنية مع مقاولة التشغيل المؤقت في أداء التعويضات المستحقة.

وقد استند القضاء في تأسيس هذا المبدأ إلى المقتضيات العامة المنصوص عليها في المواد \textenglish{35} و\textenglish{36} من مدونة الشغل المتعلقة بمنع التعسف في إنهاء علاقة الشغل، وكذا إلى مقتضيات المادة \textenglish{501} التي تُحدد المدة القصوى لعقد المهمة وتُلزم الأطراف باحترام الآجال المتفق عليها. ويرى الفقه والقضاء أن الأجير المؤقت، رغم هشاشة وضعيته القانونية، يتمتع بحد أدنى من الحماية ضد الإنهاء التعسفي خلال مدة مهمته، وهو ما يُجسد الطابع الحمائي لقانون الشغل الذي لا يميز في ضماناته الأساسية بين الأجير الدائم والأجير المؤقت.

كما أن إنهاء المهمة المبكر يُلحق بالأجير ضرراً مضاعفاً، إذ يفقده ليس فقط أجره عن المدة المتبقية، بل أيضاً فرصته في الحصول على تعويض نهاية المهمة المحتسب على أساس مجموع الأجور، فضلاً عن حرمانه من إمكانية اكتساب الخبرة والأقدمية التي قد تُؤهله للتثبيت الدائم.

\subsection{قرارات قضائية مرجعية لمحكمة النقض}
تجسدت المبادئ السابقة في أحكام وقرارات قضائية بارزة شكلت اجتهاداً مستقراً في ساحات القضاء المغربي:

\begin{tcolorbox}[colback=lightgray, colframe=maincolor, arc=5pt,
    title={\textbf{قرار محكمة النقض رقم 981 بتاريخ 18 غشت 2011 (ملف اجتماعي رقم 1691/5/1/2010)}}]
\textbf{القاعدة القضائية:}
\begin{quote}
\textit{«حيث إن لجوء المستعمل إلى أجراء مقاولة التشغيل المؤقت لا يجعل منه مشغلاً بالمفهوم القانوني، وإنما تبقى علاقة التبعية قائمة بين مقاولة التشغيل المؤقت والأجير الذي وضع رهن إشارة المقاولة المستعملة، ما دامت مقاولة التشغيل المؤقت هي التي تشغل الأجراء وتؤدي لهم أجورهم وتفي بكل الالتزامات القانونية الناشئة عن تنفيذ عقد الشغل أو إنهائه... إلا أن هذا المبدأ مشروط باحترام مقاولة التشغيل المؤقت والمقاولة المستعملة للضوابط والشروط القانونية المنصوص عليها في المواد \textenglish{495} وما يليها من مدونة الشغل، وعدم استخدام العقد للتهرب من أحكام العقود غير المحددة المدة.»}
\end{quote}
\end{tcolorbox}

\begin{tcolorbox}[colback=lightgray, colframe=sectioncolor, arc=5pt,
    title={\textbf{قرار محكمة النقض رقم 1142 بتاريخ 24 أكتوبر 2013 (ملف اجتماعي رقم 125/5/1/2013)}}]
\textbf{إعادة تكييف العقد عند سد حاجيات هيكلية دائمة:}
تتلخص وقائع القضية في مطالبة أجير تم تشغيله لعدة سنوات عبر وكالة وساطة في العمل داخل وحدة صناعية، حيث اعتبرت المحكمة أن استمرار الأجير في ممارسة مهام ترتبط مباشرة بالنشاط الإنتاجي الدائم للمؤسسة، وخضوعه لإشراف مهندسي الشركة المستعملة، يخرج العلاقة عن نطاق العمل المؤقت المنصوص عليه في المادة \textenglish{496} من مدونة الشغل.
\vspace{0.2cm}\\
وقضت محكمة النقض بتأييد الحكم القاضي بـ\textbf{إعادة تكييف العقد} واعتباره عقداً غير محدد المدة رابطاً بين الأجير والشركة المستعملة، وإلزام هذه الأخيرة بأداء كافة التعويضات المقررة عن الفصل التعسفي (التعويض عن الإشعار، ومكافأة نهاية الخدمة، والتعويض عن الضرر).
\end{tcolorbox}

\begin{tcolorbox}[colback=lightgray, colframe=accentcolor, arc=5pt,
    title={\textbf{قرار محكمة النقض رقم 456 بتاريخ 7 مارس 2016 (ملف اجتماعي رقم 832/5/1/2015)}}]
\textbf{المسؤولية التضامنية للمقاولة المستعملة عن حقوق الضمان الاجتماعي:}
تتعلق وقائع هذه القضية بأجير مؤقت اشتغل لمدة تجاوزت ثلاث سنوات داخل مقاولة صناعية عبر شركة تشغيل مؤقت، وعند مطالبته بالاستفادة من خدمات الصندوق الوطني للضمان الاجتماعي (\textenglish{CNSS})، تبين أن مقاولة التشغيل المؤقت لم تقم بالتصريح بأجوره ولم تسدد اشتراكاته طيلة فترة عمله.
\vspace{0.2cm}\\
وقد قضت محكمة النقض بأن \textbf{المقاولة المستعملة تتحمل المسؤولية التضامنية} مع مقاولة التشغيل المؤقت فيما يخص أداء اشتراكات الضمان الاجتماعي والتعويضات المترتبة عن عدم التصريح، وذلك بالنظر إلى أن المقاولة المستعملة استفادت فعلياً من خدمات الأجير واستعملت قوة عمله في نشاطها الإنتاجي دون أن تتحقق من احترام مقاولة التشغيل المؤقت لالتزاماتها القانونية تجاه مؤسسات الحماية الاجتماعية. كما أكدت المحكمة أن إخلال مقاولة التشغيل المؤقت بواجب التصريح لا يُعفي المقاولة المستعملة من مسؤوليتها، تطبيقاً لروح المواد \textenglish{89} و\textenglish{90} من مدونة الشغل.
\end{tcolorbox}

%% ─────────────────────────────────────────────────────────
%%  ثانياً: موقف القضاء من المناولة المقنعة
%% ─────────────────────────────────────────────────────────
\section{موقف القضاء من المناولة المقنعة}

تصدى القضاء المغربي بحزم لظاهرة «المناولة المقنعة» (\textenglish{Fausse sous-traitance / Prêt illicite de main-d'œuvre})، وهي الصورة التي تعمد فيها شركة إلى التعاقد مع مقاولة صورية بموجب عقد ظاهره «مناولة خدمات أو أشغال»، بينما جوهره الحقيقي يقتصر على مجرد «توريد ووضع عمال رهن الإشارة» لتأدية المهام الاعتيادية للمقاولة الأصلية، بأجور متدنية وبدون حقوق الأقدمية.

\subsection{قرائن كشف المناولة المقنعة وتوريد اليد العاملة غير المشروع}
يستند القضاء المغربي في كشف المناولة المقنعة إلى حزمة من القرائن الواقعية الدامغة:
\begin{enumerate}[leftmargin=*, itemsep=6pt]
    \item \textbf{طبيعة المهام المؤداة:} ممارسة عمال المناولة لنفس المهام التي يمارسها الأجراء الدائمون بالمقاولة الأصلية جنباً إلى جنب وفي نفس خط الإنتاج.
    \item \textbf{ملكية وسائل الإنتاج والعتاد:} استخدام أجراء المناولة لأدوات، وآلات، ومواد خام، ومقرات عمل تملكها المقاولة الأصلية حصراً، دون أي مساهمة مادية من شركة المناولة.
    \item \textbf{غياب الإشراف التقني المستقل:} انعدام أي إشراف فني أو تقني من جانب مسؤولي شركة المناولة داخل فضاء العمل، وتلقي الأجراء الأوامر والتعليمات اليومية وجداول العمل مباشرة من أطر المقاولة الأصلية.
    \item \textbf{انعدام الاستقلال المالي والفني:} عدم تمتع شركة المناولة بأي استقلال مالي أو تخصص مهني حقيقي، واقتصار دورها على دور «وسيط محاسبي» لتوزيع الأجور.
    \item \textbf{استمرارية العلاقة التعاقدية لفترات طويلة:} يُعد تجديد عقود المناولة بشكل متواصل لسنوات عديدة قرينة قوية على انتفاء الطابع المؤقت أو العرضي للخدمة المقدمة؛ فالمناولة الحقيقية تفترض بطبيعتها تعاقداً محدداً في الزمن يرتبط بإنجاز مشروع أو مهمة معينة، بينما استمرار العلاقة لأمد غير محدد يكشف عن أن الهدف الفعلي ليس الاستفادة من خبرة تقنية متخصصة، وإنما مجرد التحايل على قواعد التشغيل المباشر لتفادي الالتزامات المترتبة عنه.
    \item \textbf{عدم وجود خبرة تقنية خاصة لشركة المناولة:} يُشكل افتقار شركة المناولة إلى تخصص مهني أو تقني مميز عن نشاط المقاولة الأصلية قرينة حاسمة على صورية عقد المناولة. فالمناولة المشروعة تقوم على إسناد أعمال تتطلب كفاءات خاصة لا تتوفر لدى المقاولة الأصلية، أما حين تكون شركة المناولة مجرد هيكل إداري لتوظيف العمال وأداء أجورهم دون أي قيمة مضافة تقنية أو فنية، فإن ذلك يُفضي حتماً إلى تكييف العقد باعتباره إعارة غير مشروعة لليد العاملة.
\end{enumerate}

\subsection{تطبيقات قضائية لمحاكم الاستئناف (الدار البيضاء والرباط)}

\begin{tcolorbox}[colback=lightgray, colframe=maincolor, arc=5pt,
    title={\textbf{قرار محكمة الاستئناف بالدار البيضاء رقم 3241 بتاريخ 12 نونبر 2018}}]
\textbf{القاعدة المستخلصة:}
أكدت الغرفة الاجتماعية بمحكمة الاستئناف بالدار البيضاء على المسؤولية المباشرة للمقاولة المستعملة عند خرق الشروط المنصوص عليها في المادتين \textenglish{496} و\textenglish{500} من مدونة الشغل؛ حيث قررت أن تشغيل الأجير المؤقت لأكثر من المدد القصوى المسموح بها قانوناً دون سند مبرر، وتكليفه بمهام دائمة لا تتصل بأي زيادة ظرفية في الإنتاج، يجعل المقاولة المستعملة مشغلاً أصلياً للأجير، وتتحمل بالتالي المسؤولية الكاملة عن طرده دون اتباع مسطرة الفصل التأديبي المنصوص عليها في المادة \textenglish{62} من المدونة.
\end{tcolorbox}

\begin{tcolorbox}[colback=lightgray, colframe=sectioncolor, arc=5pt,
    title={\textbf{قرار محكمة الاستئناف بالرباط رقم 789 بتاريخ 16 أكتوبر 2019}}]
\textbf{تطبيق المسؤولية التضامنية (المادتان 89 و90 من مدونة الشغل):}
قضت المحكمة بإلزام الشركة الأصلية تضامنياً مع مقاول الباطن بأداء مستحقات الأجراء المتخلفة عن الأجور والتعويض عن العطلة السنوية، بعدما ثبت عجز مقاول الباطن وتملصه من أداء الواجبات، مكرسةً حق الأجراء في ممارسة \textbf{الدعوى المباشرة} ضد الشركة صاحبة المشروع الأصلي بصريح نص المادة \textenglish{90}.
\end{tcolorbox}

\begin{tcolorbox}[colback=lightgray, colframe=accentcolor, arc=5pt,
    title={\textbf{حكم المحكمة الابتدائية بالدار البيضاء رقم 2156 بتاريخ 23 يناير 2020}}]
\textbf{تكييف عقد النظافة باعتباره توريداً مقنعاً لليد العاملة:}
تتعلق وقائع هذه القضية بشركة مقاولات أبرمت عقداً مع شركة متخصصة ظاهرياً في «خدمات النظافة والصيانة»، غير أن التحقيق الميداني كشف أن هذه الشركة لا تمتلك أي معدات نظافة خاصة بها ولا مواد تنظيف، وأن عمالها يستعملون حصرياً أدوات ومواد تملكها المقاولة الأصلية. كما تبين أن هؤلاء العمال لا يقتصرون على مهام النظافة، بل يُكلفون أيضاً بأعمال الشحن والتفريغ ونقل البضائع داخل المخازن، وهي مهام تدخل في صلب النشاط العادي والمعتاد للمقاولة الأصلية.
\vspace{0.2cm}\\
وقد قضت المحكمة بأن العقد المبرم بين الطرفين \textbf{لا يُعد عقد مناولة حقيقياً}، وإنما يُشكل في حقيقته توريداً مقنعاً لليد العاملة يتعارض مع أحكام المواد \textenglish{86} إلى \textenglish{90} من مدونة الشغل، وأعادت تكييف العلاقة باعتبار المقاولة الأصلية هي المشغل الحقيقي لجميع العمال، مع إلزامها بأداء كافة حقوقهم المكتسبة منذ بداية اشتغالهم الفعلي لديها، بما في ذلك علاوات الأقدمية ومكافآت نهاية الخدمة.
\end{tcolorbox}

%% ─────────────────────────────────────────────────────────
%%  ثالثاً: التجربة المقارنة
%% ─────────────────────────────────────────────────────────
\section{التجربة المقارنة}

يُبرز استقراء القوانين الأجنبية والعربية أن مشكلة التوفيق بين مرونة التشغيل وحماية حقوق العمال تمثل هاجساً تشريعياً كونياً، غير أن الحلول المعتمدة تتفاوت في درجات صرامتها وضماناتها:

\subsection{التجربة الفرنسية: الصرامة الزجرية وتجريم الاتجار باليد العاملة}
يُعتبر قانون العمل الفرنسي (\textenglish{Code du travail}) من أكثر الأنظمة دقة في تنظيم العمل المؤقت (\textenglish{Travail temporaire / Intérim})، حيث وضع حواجز صارمة لمنع استغلال العمال:
\begin{itemize}[leftmargin=*, itemsep=6pt]
    \item \textbf{قاعدة حظر شغل المناصب الدائمة (المادة \textenglish{L. 1251-5}):} تنص صراحة على أنه لا يجوز اللجوء إلى عقد الشغل المؤقت لشغل منصب يرتبط بالنشاط الطبيعي والدائم للمؤسسة المستعملة.
    \item \textbf{إعادة التكييف التلقائي بقوة القانون (المادة \textenglish{L. 1251-40}):} يحق للأجير المؤقت اللجوء إلى قضاء العمل (\textenglish{Conseil de prud'hommes}) للمطالبة بإعادة تكييف عقده إلى عقد غير محدد المدة (\textenglish{CDI}) مباشرة في مواجهة المقاولة المستعملة عند خرق قواعد اللجوء أو المدد، مع استحقاقه تعويضاً إجبارياً عن إعادة التكييف (\textenglish{Indemnité de requalification}) لا يقل عن أجر شهر كامل.
    \item \textbf{التجريم الجنائي للإعارة غير المشروعة للعمالة والاتجار بها:} يُعاقب القانون الفرنسي على «جريمة الاتجار باليد العاملة» (\textenglish{Délit de marchandage - Art. L. 8231-1}) و«الإعارة غير المشروعة لليد العاملة» (\textenglish{Prêt illicite de main-d'œuvre - Art. L. 8241-1}) بعقوبات حبسية تصل إلى \textenglish{3} سنوات وغرامات مالية ضخمة تصل إلى \textenglish{45.000} يورو للأشخاص الذاتيين وتتضاعف \textenglish{5} مرات للأشخاص الاعتباريين.
    \item \textbf{تعويض نهاية المهمة (\textenglish{Indemnité de fin de mission}):} إلزام المشغل بصرف تعويض نسبته \textenglish{10\%} من مجموع الأجور الإجمالية التي تقاضاها الأجير المؤقت كتعويض عن الهشاشة وعدم استقرار العمل.
    \item \textbf{مبدأ المساواة في الأجر:} إلزام المقاولة المستعملة بمنح الأجير المؤقت نفس الأجر والامتيازات التي يستفيد منها الأجير الدائم ذو المؤهلات المماثلة.
\end{itemize}

\subsection{التجربة الإسبانية: إصلاح 2021 ومحاصرة الهشاشة التعاقدية}
تعتبر التجربة الإسبانية نموذجاً حديثاً ملهماً في كبح جماح العمل المؤقت؛ فبعد عقود من انتشار العقود الهشة، أقرت إسبانيا تعديلاً جوهرياً بموجب المرسوم الملكي بقانون رقم \textenglish{32/2021} (\textenglish{Reforma Laboral 2021}):
\begin{itemize}[leftmargin=*, itemsep=5pt]
    \item افتراض أن جميع عقود الشغل مبرمة لأجل غير مسمى (\textenglish{Contrato indefinido}) كأصل مطلق، وتقييد العقود المؤقتة في أضيق الحدود (الظروف الإنتاجية الطارئة فقط وبسقف زمني صارم).
    \item فرض المساواة المطلقة في الأجور، وظروف الصحة والسلامة، والترقيات والتدريب بين أجراء شركات التشغيل المؤقت (\textenglish{Empresas de Trabajo Temporal - ETT}) والأجراء الدائمين في المقاولة المستعملة.
    \item تشديد العقوبات الإدارية والمالية على مقاولات المناولة التي يثبت تحايلها، وربط المسؤولية التضامنية بجميع الالتزامات المهنية.
\end{itemize}

\subsection{التجربة الألمانية: نموذج المشاركة والحوار الاجتماعي}
يُقدم النظام القانوني الألماني نموذجاً متميزاً في تنظيم التشغيل المؤقت يقوم على ركيزتين أساسيتين: المشاركة المؤسساتية للعمال والحوار الاجتماعي المنتظم. فقد أخضع المشرع الألماني نشاط مقاولات التشغيل المؤقت لقانون تنظيم إعارة العمال (\textenglish{Arbeitnehmerüberlassungsgesetz - AÜG}) الذي خضع لإصلاح جذري سنة \textenglish{2017} تضمن مقتضيات حمائية متقدمة:

\begin{itemize}[leftmargin=*, itemsep=6pt]
    \item \textbf{مبدأ المساواة في المعاملة من اليوم الأول (\textenglish{Equal Pay / Equal Treatment}):} أقر الإصلاح حق الأجير المؤقت في التمتع بنفس الأجر وظروف العمل المطبقة على الأجراء الدائمين في المقاولة المستعملة ابتداءً من أول يوم عمل، مع إمكانية الاستثناء بموجب اتفاقيات جماعية قطاعية لمدة أقصاها \textenglish{9} أشهر، على أن يتم التدرج في المساواة الأجرية بعد ذلك وصولاً إلى المساواة الكاملة بعد \textenglish{15} شهراً كحد أقصى.
    \item \textbf{سقف زمني أقصى للمهمة (\textenglish{18} شهراً):} حدد القانون مدة قصوى لوضع الأجير المؤقت رهن إشارة نفس المقاولة المستعملة بـ\textenglish{18} شهراً متواصلة، وعند تجاوز هذه المدة يُعاد تكييف العقد تلقائياً إلى عقد غير محدد المدة يربط الأجير بالمقاولة المستعملة مباشرة.
    \item \textbf{دور مجالس العمل (\textenglish{Betriebsrat}):} يتمتع مجلس العمل (مجلس المقاولة) بصلاحيات واسعة في مراقبة اللجوء إلى التشغيل المؤقت داخل المؤسسة، إذ يحق له الاطلاع على جميع عقود الوضع رهن الإشارة، والاعتراض على أي تشغيل مؤقت يعتبره مخالفاً للقانون أو ماساً بحقوق الأجراء الدائمين، مما يُشكل رقابة داخلية فعالة ومباشرة تُكمل الرقابة القضائية اللاحقة.
    \item \textbf{الاتفاقيات الجماعية القطاعية (\textenglish{Tarifverträge}):} يلعب التفاوض الجماعي على مستوى القطاعات الصناعية دوراً محورياً في تحسين أوضاع العمال المؤقتين بألمانيا؛ إذ تُبرم اتفاقيات جماعية بين نقابات العمال واتحادات أصحاب المقاولات تتضمن شروطاً أكثر ملاءمة من الحد الأدنى القانوني، لا سيما في مجالات العلاوات الخاصة بالقطاعات الصناعية والتكوين المهني المستمر وضمانات التثبيت بعد انتهاء المهمة.
\end{itemize}

ويُمثل النموذج الألماني مقاربة تشاركية تجمع بين المرونة التي يحتاجها الاقتصاد والحماية الفعلية للعمال، وهو ما يمكن أن يُشكل مصدر إلهام للمشرع المغربي في تطوير منظومته القانونية المتعلقة بالتشغيل المؤقت والمناولة.

\subsection{تجارب عربية مقارنة (تونس ومصر)}
\begin{itemize}[leftmargin=*, itemsep=6pt]
    \item \textbf{التجربة التونسية:} أفردت «مجلة الشغل التونسية» في المواد \textenglish{28} وما يليها تنظيماً حذراً لمقاولي اليد العاملة، حيث منعت المتاجرة بالأجراء واشترطت موافقة السلطة الإدارية. وتعيش تونس اليوم تحولاً تشريعياً ومجتمعياً نوعياً نحو المنع الشامل لظاهرة المناولة في الوظيفة العمومية والقطاع الخاص، واعتبارها شكلاً من أشكال العمل الهش الذي يتعارض مع مبدأ العدالة الاجتماعية وكرامة الشغالين.
    \item \textbf{التجربة المصرية:} نظّم قانون العمل المصري رقم \textenglish{12} لسنة \textenglish{2003} مسألة تشغيل العمال عبر متعهدي ووكالات التشغيل في الداخل والخارج؛ وحظر تشغيل العمال عن طريق وسطاء غير مرخص لهم، كما منع تحصيل أي مبالغ مالية من الأجراء مقابل إلحاقهم بالعمل، وألزم صاحب العمل الأصلي بالتضامن الكامل في الوفاء بالأجور والمستحقات العمالية عند الاستعانة بمقاولي الباطن.
\end{itemize}

%% ─────────────────────────────────────────────────────────
%%  رابعاً: المعايير الدولية للشغل
%% ─────────────────────────────────────────────────────────
\section{المعايير الدولية للشغل}

تُشكل معايير منظمة العمل الدولية (\textenglish{OIT}) الإطار المرجعي الأممي الذي يوجه السياسات الوطنية في تنظيم أسواق العمل:

\begin{enumerate}[leftmargin=*, itemsep=8pt]
    \item \textbf{الاتفاقية الدولية رقم 181 (لسنة 1997) بشأن وكالات الاستخدام الخاصة:}
          تعترف هذه الاتفاقية بدور وكالات التشغيل المؤقت في مرونة سوق الشغل، لكنها تُخضع نشاطها لمنظومة مشددة من الضمانات:
          \begin{itemize}[leftmargin=*, itemsep=3pt]
              \item منع وكالات التشغيل من اقتطاع أي أتعاب أو مصاريف مباشرة أو غير مباشرة من أجور العمال.
              \item ضمان الحق في المفاوضة الجماعية والحرية النقابية لعمال التشغيل المؤقت.
              \item إلزام الدول الأطراف بوضع آليات لتحديد مسؤوليات كل من الوكالة والمقاولة المستعملة في مجالات الأجور وحوادث الشغل.
          \end{itemize}
    \item \textbf{التوصية رقم 188 (لسنة 1997):} تدعو الدول إلى وضع آليات رقابية فعالة وعقوبات رادعة (بما فيها شطب وسحب تراخيص المقاولات المخالفة).
    \item \textbf{الاتفاقية رقم 94 بشأن بنود العمل في العقود العامة:} إلزام المقاولين وأصحاب الصفقات باحترام شروط الشغل اللائقة للأجراء المناولين.
    \item \textbf{مفهوم العمل اللائق (\textenglish{Decent Work}):} يؤكد إعلان منظمة العمل الدولية بشأن العدالة الاجتماعية من أجل عولمة منصفة (\textenglish{2008}) أن العمل ليس سلعة تجارية تباع وتشترى، وأن المرونة الاقتصادية تفقد مشروعيتها متى أفضت إلى حرمان الأجراء من الأمان الوظيفي والحماية الاجتماعية.
    \item \textbf{الاتفاقية رقم 122 (لسنة 1964) بشأن سياسة الاستخدام:}
          تُلزم هذه الاتفاقية الدول المصادقة عليها بانتهاج سياسة وطنية نشيطة ترمي إلى تحقيق \textbf{الاستخدام الكامل والمنتج والمختار بحرية}، بما يضمن لكل قادر على العمل وباحث عنه فرصة الحصول على شغل لائق ومستقر. وتُشكل هذه الاتفاقية أساساً معيارياً مهماً لمساءلة الدول التي تتوسع في اعتماد أشكال التشغيل الهش والمؤقت على حساب الاستخدام المستقر؛ إذ ينبغي ألا يُوظف التشغيل المؤقت كبديل منهجي عن خلق مناصب شغل دائمة ومنتجة، بل يجب أن يظل استثناءً محصوراً في الحاجيات الظرفية الحقيقية للمقاولات.
    \item \textbf{أجندة العمل اللائق \textenglish{2030} وأهداف التنمية المستدامة:}
          ربطت منظمة العمل الدولية بين مفهوم العمل اللائق وأهداف التنمية المستدامة (\textenglish{SDGs})، لا سيما \textbf{الهدف الثامن} (\textenglish{SDG 8}) المتعلق بـ«تعزيز النمو الاقتصادي المطرد والشامل والمستدام، والعمالة الكاملة والمنتجة، وتوفير العمل اللائق للجميع». ويستلزم هذا الهدف من الدول الأعضاء اتخاذ تدابير فعالة للقضاء على العمل الجبري والاتجار بالبشر وعمل الأطفال، وضمان بيئة عمل آمنة ومأمونة لجميع العمال بمن فيهم العمال المهاجرون والمؤقتون. كما تُلزم الغاية \textenglish{8.8} من هذا الهدف الدول بحماية حقوق العمل وتعزيز بيئات العمل المأمونة لجميع العمال، بما فيهم العمال الذين يعملون في ظروف هشة عبر عقود التشغيل المؤقت والمناولة.
\end{enumerate}

\subsection{تقييم مدى ملاءمة التشريع المغربي للمعايير الدولية}
يكشف استقراء المعايير الدولية المذكورة ومقارنتها بالمنظومة التشريعية المغربية عن وجود فجوات تستدعي التدخل التشريعي لسدها وتحقيق الملاءمة المطلوبة:

\begin{itemize}[leftmargin=*, itemsep=6pt]
    \item \textbf{على مستوى الضمانات الأجرية:} رغم نص المادة \textenglish{504} من مدونة الشغل على مبدأ المساواة في الأجر، إلا أن المشرع المغربي لم ينص على تعويض إلزامي عن الهشاشة (\textenglish{Précarité}) على غرار ما هو معمول به في فرنسا (نسبة \textenglish{10\%} من مجموع الأجور)، مما يُضعف الحماية المالية للأجير المؤقت ولا يُعوضه عن عدم استقرار وضعيته المهنية.
    \item \textbf{على مستوى الجزاءات:} تبقى الغرامات المنصوص عليها في مدونة الشغل المغربية ضعيفة وغير رادعة مقارنة بالعقوبات المقررة في التشريعات المتقدمة، كما يغيب التجريم الصريح لجريمة «الاتجار باليد العاملة» و«الإعارة غير المشروعة» بوصفهما جنحتين مستقلتين تستوجبان عقوبات حبسية.
    \item \textbf{على مستوى الرقابة المؤسساتية:} لم يُفرد المشرع المغربي دوراً مؤسساتياً لممثلي الأجراء أو المؤسسات النقابية في مراقبة اللجوء إلى التشغيل المؤقت والمناولة داخل المقاولات، على خلاف النموذج الألماني الذي يمنح مجالس العمل صلاحيات رقابية واسعة في هذا الشأن.
    \item \textbf{على مستوى المصادقة على الاتفاقيات:} لم يُصادق المغرب بعد على الاتفاقية الدولية رقم \textenglish{181} بشأن وكالات الاستخدام الخاصة، مما يُفقد الأجراء المؤقتين مرجعية دولية ملزمة يمكن الاحتجاج بها أمام القضاء الوطني لتعزيز حمايتهم.
\end{itemize}

%% ─── خلاصة الفصل ──────────────────────────────────────────
\section*{خلاصة الفصل}
\addcontentsline{toc}{section}{خلاصة الفصل}

يتبين من خلال استقراء الاجتهاد القضائي المغربي والتجارب المقارنة أن القضاء لم يقف مكتوف الأيدي أمام مخاطر استغلال الشغل المؤقت والمناولة كأدوات للتهرب من تكاليف قانون الشغل، بل سارع—مستعيناً بسلطته في تكييف العقود وبحث عناصر التبعية الواقعية—إلى ردع الممارسات الصورية وإعادة الأجراء إلى مظلة العقود الدائمة والمستقرة.

غير أن هذا المجهود القضائي، على عظمته، يظل علاجاً لاحقاً يقتضي من الأجير خوض غمار مساطر تقاضٍ طويلة ومكلفة في ظل تفاوت موازين القوى مع المقاولات الكبرى. وبالمقارنة مع التشريعات الدولية المتقدمة (كالنموذج الفرنسي والإسباني والألماني)، يتضح أن التشريع المغربي بحاجة ملحة إلى التطوير وسد الثغرات النصية، لا سيما بالنص الصريح على جزاء إعادة التكييف ورفع الغرامات وتجريم الإعارة غير المشروعة لليد العاملة كجنحة مستقلة.

ويستدعي الوضع الراهن تدخلاً تشريعياً عاجلاً يستلهم أفضل الممارسات الدولية، ولا سيما إقرار تعويض إلزامي عن الهشاشة لفائدة الأجراء المؤقتين، وتعزيز آليات الرقابة الإدارية والنقابية على مقاولات التشغيل المؤقت والمناولة، فضلاً عن المصادقة على الاتفاقيات الدولية ذات الصلة. إن الاجتهاد القضائي الفاعل، رغم أهميته البالغة في سد الفراغ التشريعي وحماية الحقوق الأساسية للعمال، لا يمكن أن يحل محل إرادة تشريعية واضحة وحازمة تُرسي قواعد آمرة لا تقبل التأويل ولا التحايل، وهو ما يشكل موضوع الدراسة والتحليل في الفصل العاشر الموالي.
